\documentclass[11pt,a4paper]{article}
\usepackage[T1]{fontenc}
\usepackage[utf8]{inputenc}
\usepackage[french]{babel}
\usepackage{lmodern}
\usepackage{amsmath,amssymb,mathtools}
\usepackage{geometry}
\geometry{top=22mm,bottom=22mm,left=23mm,right=23mm,headheight=15pt,headsep=6mm}
\usepackage{microtype,enumitem,booktabs,array,fancyhdr,lastpage,needspace}
\usepackage[hidelinks]{hyperref}
\setlength{\parindent}{0pt}
\setlength{\parskip}{6pt}
\setlist[enumerate]{leftmargin=*,label=\textbf{\alph*)},itemsep=8pt,topsep=4pt}
\pagestyle{fancy}\fancyhf{}
\renewcommand{\headrulewidth}{0.3pt}
\fancyfoot[C]{\small\thepage\,/\,\pageref{LastPage}}
\fancyfoot[L]{\scriptsize Version de travail -- 6 septembre 2026}
\setlength{\emergencystretch}{2em}
\hypersetup{pdftitle={MAT0339 -- Série B, examen 2 -- sujet},pdfauthor={Document de travail pédagogique}}
\fancyhead[L]{\small MAT0339 -- Série B -- Examen 2}
\fancyhead[R]{\small SUJET À VALIDER}
\begin{document}
{\LARGE\bfseries Série B -- Examen 2}\par
{\large Dénombrement, probabilités et algèbre linéaire}\par
\textbf{Épreuve étudiante -- version de travail}\par
Portée : chapitres 7 à 12. Total : 100 points. Durée proposée : 120 min.\par
\textbf{Nom :} \hrulefill\quad \textbf{Groupe :} \rule{22mm}{0.3pt}\par
\small Montrer les démarches. Conserver les valeurs exactes jusqu’à l’arrondi final; annoncer les unités. Convention provisoire : calculatrice scientifique non symbolique autorisée, sans notes. Les modalités doivent être confirmées par l’enseignant. Les pages suivantes offrent une page par question; une feuille supplémentaire peut être jointe.\normalsize\par
\vspace{3mm}\hrule\vspace{4mm}
{\large\bfseries Question 1 -- Compter les issues}\hfill\textbf{15 points}\par
\begin{enumerate}
\item \textbf{(5 points)} Un code comprend une lettre parmi les 26 lettres, suivie de trois chiffres distincts. Le zéro est permis à toute position numérique. Combien de codes existent? Justifier la méthode.
\item \textbf{(5 points)} Combien de comités sans rôle de quatre personnes peut-on choisir parmi neuf personnes distinctes? Justifier la méthode.
\item \textbf{(5 points)} Combien d’anagrammes distinctes peut-on former avec LEVEL? Toutes les lettres sont utilisées. Justifier la méthode.
\end{enumerate}
\vspace{3mm}\textit{Démarche et réponse :}\par\vfill

\newpage
{\large\bfseries Question 2 -- Tirages sans remise}\hfill\textbf{20 points}\par
\begin{enumerate}
\item \textbf{(8 points)} Une urne contient 5 boules rouges et 3 boules bleues. On tire uniformément deux boules successivement, sans remise. On note \(R_i\) et \(B_i\) la couleur au tirage \(i\). Construire un arbre pondéré et calculer les probabilités des quatre chemins.
\item \textbf{(4 points)} Calculer la probabilité d’obtenir deux boules de couleurs différentes.
\item \textbf{(4 points)} Calculer \(P(B_1\mid R_2)\).
\item \textbf{(4 points)} Les événements \(R_1\) et \(R_2\) sont-ils indépendants? Justifier par une égalité ou une inégalité de probabilités.
\end{enumerate}
\vspace{3mm}\textit{Démarche et réponse :}\par\vfill

\newpage
{\large\bfseries Question 3 -- Essais répétés et espérance}\hfill\textbf{15 points}\par
\begin{enumerate}
\item \textbf{(4 points)} On effectue 4 essais indépendants, chacun de probabilité de succès \(p=0{,}75\). Pour \(X\), le nombre de succès, identifier le modèle et justifier les hypothèses.
\item \textbf{(4 points)} Calculer \(P(X=3)\).
\item \textbf{(3 points)} Calculer la probabilité d’au moins un succès.
\item \textbf{(4 points)} Pour un seul essai, le gain net est de 5~\$ en cas de succès et une perte de 3~\$ sinon. Calculer et interpréter l’espérance du gain net.
\end{enumerate}
\vspace{3mm}\textit{Démarche et réponse :}\par\vfill

\newpage
{\large\bfseries Question 4 -- Angle et projection}\hfill\textbf{15 points}\par
\begin{enumerate}
\item \textbf{(3 points)} Soient \(\mathbf u=(1,3),\quad\mathbf v=(2,2)\). Calculer les deux normes.
\item \textbf{(5 points)} Calculer le produit scalaire et l’angle entre les vecteurs, au centième de degré.
\item \textbf{(7 points)} Calculer \(\operatorname{proj}_{\mathbf v}\mathbf u\), puis vérifier que le vecteur résiduel est orthogonal à \(\mathbf v\).
\end{enumerate}
\vspace{3mm}\textit{Démarche et réponse :}\par\vfill

\newpage
{\large\bfseries Question 5 -- Composer et inverser des matrices}\hfill\textbf{20 points}\par
\begin{enumerate}
\item \textbf{(6 points)} Soient \[A=\begin{pmatrix}1 & 2\\2 & 3\end{pmatrix},\qquad B=\begin{pmatrix}2 & 0\\0 & 1\end{pmatrix}.\] Calculer \(AB\) et \(BA\). Comparer les résultats.
\item \textbf{(6 points)} Calculer \(\det A\), puis \(A^{-1}\). Vérifier le produit \(AA^{-1}\).
\item \textbf{(8 points)} Résoudre \(A\mathbf x=\begin{pmatrix}8\\13\end{pmatrix}\). Vérifier dans les deux équations initiales.
\end{enumerate}
\vspace{3mm}\textit{Démarche et réponse :}\par\vfill

\newpage
{\large\bfseries Question 6 -- Intersection et classification d’un système}\hfill\textbf{15 points}\par
\begin{enumerate}
\item \textbf{(8 points)} La droite \((x,y,z)=(0,1,1)+t(1,2,-1)\), \(t\in\mathbb R\), rencontre-t-elle le plan \(2x-y+z=3\)? Déterminer leur intersection et justifier son unicité.
\item \textbf{(7 points)} En fonction du réel \(k\), classer les solutions du système \[\begin{cases}x+y+z=4,\\x-y+z=2,\\2x+2z=k.\end{cases}\] Donner une paramétrisation lorsqu’il est compatible.
\end{enumerate}
\vspace{3mm}\textit{Démarche et réponse :}\par\vfill
\end{document}
